\documentclass[12pt]{article}
\usepackage[T1]{fontenc}
\usepackage[utf8]{inputenc}
\usepackage{lmodern}
\usepackage{textcomp}
\usepackage{enumitem}
\usepackage{geometry}
\geometry{margin=1in}
\begin{document}
\begin{center}
Answer Key - Astronomy - Graduate Level Assessment 21 \
\small (Answers are ordered exactly as in the question paper.)
\end{center}

\section*{Multiple Choice}
\begin{enumerate}[label=\arabic*.]
\item \textbf{Answer:} D
\\ \textit{Options:} A) Power emitted is 1/16 times as high; peak emission wavelength is 1/2 as long.; B) Power emitted is 1/4 times as high; peak emission wavelength is 2 times longer.; C) Power emitted is 1/4 times as high; peak emission wavelength is 1/2 as long.; D) Power emitted is 1/16 times as high; peak emission wavelength is 2 times longer.
\\ \textit{Note:} According to Wien's displacement law, as the temperature of a blackbody halved, the peak emission wavelength doubles, and according to the Stefan-Boltzmann law, the total power emitted decreases by a factor of 1/16.
\item \textbf{Answer:} C
\\ \textit{Options:} A) Mass 5 kg weight 100 N; B) Mass 5 kg weight 50 N; C) Mass 5 kg weight 25 N; D) Mass 10 kg weight 100 N
\\ \textit{Note:} The mass of the sandbag remains constant at 5 kg regardless of the gravitational field, while its weight is halved to 25 N on a planet with half the gravity compared to Earth, since weight is calculated as mass multiplied by the gravitational acceleration.
\item \textbf{Answer:} A
\\ \textit{Options:} A) one medium sized satellite and two small satellites.; B) no satellites.; C) one large satellite and three small satellites.; D) one large satellite.
\\ \textit{Note:} The correct answer is supported by recent astronomical observations that confirm Pluto has one medium-sized moon, Charon, and two smaller moons, Nix and Hydra, positioning it as a complex system within our solar system.
\item \textbf{Answer:} D
\\ \textit{Options:} A) Mars does not have enough internal heat to drive the greenhouse effect; B) Mars is too far from the sun for the greenhouse effect to work; C) the greenhouse effect requires an ozone layer which Mars does not have; D) the atmosphere on Mars is too thin to trap a significant amount of heat
\\ \textit{Note:} The correct answer is based on the fact that Mars' thin atmosphere, with a surface pressure less than 1\% of Earth's, is insufficient to retain heat effectively, despite its high carbon dioxide content.
\item \textbf{Answer:} D
\\ \textit{Options:} A) Will the Sun become a black hole?; B) Is the universe infinitely large?; C) How old is the visible universe?; D) Are we alone in the universe?
\\ \textit{Note:} The Drake equation is designed to estimate the number of active, communicative extraterrestrial civilizations in the Milky Way galaxy, thus directly addressing the question of whether humanity is the only intelligent life in the universe.
\end{enumerate}

\section*{True / False}
\begin{enumerate}[label=\arabic*.]
\setcounter{enumi}{5}
\item \textbf{Answer:} False
\\ \textit{Options:} A) True; B) False
\\ \textit{Note:} A black body is an idealized physical object that absorbs all incident electromagnetic radiation without reflecting any, which is the defining characteristic that makes the statement false.
\item \textbf{Answer:} True
\\ \textit{Options:} A) True; B) False
\\ \textit{Note:} The statement is true because the Phoenix Mars Lander was indeed designed to land in the Martian polar region on May 25, 2008, with the specific mission objective of studying the presence and characteristics of ice-rich soil.
\item \textbf{Answer:} False
\\ \textit{Options:} A) True; B) False
\\ \textit{Note:} A scientific theory is never conclusively proven correct based on a single experiment because scientific theories are built on a body of evidence from multiple experiments and observations that must consistently support the theory across varying conditions and contexts.
\item \textbf{Answer:} False
\\ \textit{Options:} A) True; B) False
\\ \textit{Note:} The Large Magellanic Cloud is a satellite galaxy of the Milky Way and contains many stars and nebulae, but it is not classified as a planetary nebula; the closest known planetary nebula to Earth is actually the Helix Nebula.
\item \textbf{Answer:} True
\\ \textit{Options:} A) True; B) False
\\ \textit{Note:} Venus does not experience seasons because its axial tilt is only about 3 degrees, which means that there is minimal variation in sunlight distribution throughout its orbit around the Sun.
\end{enumerate}

\section*{Long Form Response}
\begin{enumerate}[label=\arabic*.]
\setcounter{enumi}{10}
\item \textbf{Answer:} Sidereal days and solar days differ primarily in their definitions: a sidereal day, approximately 23 hours, 56 minutes, is the time it takes for the Earth to complete one full rotation relative to the stars, while a solar day, about 24 hours, is based on the Earth's rotation relative to the Sun. The statement that the Moon's circuit in our sky corresponds to one solar day is misleading; although the Moon appears to move through the sky over the course of a solar day, it actually takes about 27.3 days to complete one full orbit around the Earth. Consequently, the Moon's position in relation to the stars changes more slowly than the daily solar cycle. This distinction is crucial for understanding the dynamics of celestial motion and reinforces the need for precise definitions in astronomical observations.
\\ \textit{Note:} correct because it accurately distinguishes between sidereal and solar days based on Earth's rotational references, and clarifies that the Moon's movement across the sky does not align with the duration of a solar day, as its actual orbital period is approximately 27.3 days.
\item \textbf{Answer:} The differentiation in material composition between the inner and outer planets of our solar system can be attributed primarily to the thermal gradient present in the solar nebula during its formation. In the hotter inner region of the nebula, only metals and silicate minerals could condense, leading to the formation of terrestrial planets characterized by rocky compositions. In contrast, the cooler outer regions allowed for the condensation of hydrogen compounds, such as water, ammonia, and methane, which are abundant but could not form in the higher temperatures of the inner nebula. This thermal gradient, therefore, resulted in the inner planets being composed predominantly of heavy, rocky materials, while the outer planets formed as gas giants with substantial icy and gaseous components. The processes of accretion and the resulting planetary formation were thus heavily influenced by temperature variations within the protoplanetary disk.
\\ \textit{Note:} The answer correctly identifies the thermal gradient in the solar nebula as the key factor influencing the composition of the inner and outer planets, with higher temperatures facilitating the condensation of only rocky materials in the inner region, while the cooler outer region permitted the formation of gas-rich planets with hydrogen compounds.
\end{enumerate}

\vfill
\noindent\textit{End of Answer Key}
\end{document}